
\def\PUDcodigo{ACE021}

\if\COMPILAR0
   \def\PUDcodacad{04507.59}
\else
   \def\PUDcodacad{04506.57}
\fi

\def\PUDdisciplina{Educação Física}

\def\PUDsemestre{Opt}

\def\PUDhoras{40}

%NOVOS ---------------------------------------------------
\def\PUDhorasTeoria{20}
\def\PUDhorasPratica{20}

\def\PUDinterECA{---}

\def\PUDatuaisECA{---}

\def\PUDrecursos{Práticas esportivas realizadas no ginásio, campo e/ou academia, com materias esportivos do campus. Eventualmente, podem ser realizadas aulas em sala de aula com projetor multimídia, quadro branco e pincel.}

%----------------------------------------------------------

\def\PUDcreditos{2}

\def\PUDprerequisito{---}

\def\PUDpreacad{---}

\def\PUDobjetivos{Compreender a importância da atividade física e dos esportes para o desenvolvimento biopsicossocial, enquanto ser consciente e comprometido com o seu contexto histórico, por meio da autonomia, ludicidade, prazer e reflexão crítica.}

\def\PUDementa{Metodologia dos Esportes Coletivos (futsal, futebol, voleibol, basquetebol e handebol); Exercício Físico e Aptidão Física relacionada à Saúde; Atividades Aquáticas e Natação; Atividades de Academia e Musculação.}

\def\PUDconteudo{ &  FUTSAL/FUTEBOL/BASQUETEBOL/HANDEBOL E VOLEIBOL.  Fundamentos, Técnica e Tática do futsal e futebol.  Vivenciar situações de jogos competitivos ou recreativos.
\\
&  NATAÇÃO E ATIVIDADES DE ACADEMIA.  Fundamentos e Técnicas dos 04 nados: crawl, peito, costas e borboleta.  Vivenciar a natação em situações de jogos competitivos ou recreativos.  Características e definições de exercício físico e aptidão física relacionada à saúde; Características e vivências de um exercício aeróbio e anaeróbio; Vivenciando exercícios de força, velocidade, agilidade, resistência e aptidão cardiorrespiratória. }

\def\PUDmetodo{As aulas serão 100 porcento práticas, por meio de dinâmicas de grupo, jogos, brincadeiras, circuitos e treinamentos físico-desportivos. Os alunos poderão escolher a modalidade que desejam participar de acordo com seus interesses e/ou aptidões físicas. }

\def\PUDavalia{A avaliação será formativa e somativa por meio do desempenho dos alunos nas aulas (frequência e participação). }

\def\PUDbibbasica{ &  Darido, Suraya Cristina. Educação Física na escola: Questões e Reflexões.  Guanabara Koogan, 2003.  104p.  ISBN: 9788527708364.
\\
 &  Manhães, Elaine. 519 atividades e jogos para esportes de quadra.  Rio de Janeiro: Sprint, 2011.  171p.  ISBN: 8573322918.
%&   DELAVIER, F. Guia de movimentos da musculação. Manole, 2ª ed. 2015. 
\\
 &  Santarem, José Maria. Musculação em todas as idades.  Editora: Manole. 2012.  200p.  ISBN: 9788520434352. }
%&   GONZALEZ, Fernando Jaime. DARIDO, Suraya Cristina. OLIVEIRA, Amauri Aparecido Bássoli de. Esportes de invasão. Eduem, 2014.

\def\PUDbibcomplementar{ &  COLEIVO DE AUTORES. Metodologia do ensino da Educação Física. 2º ed.  Cortez, 2009.  200p.  ISBN: 9788524915413.
\\
&  SILVA, Pedro Antonio da. 3000 exercícios e jogos para a Educação Física escolar.  v. 2.  Sprint, 2005.  ISBN: 8573321768.	
\\ 
 &  SOEIRO. Maria Isaura Plácido. SILVA, Maria Ione da. Educação Física escolar: pesquisas e reflexões.  UERN, 2014.  187p.  ISBN: 9788576210849. 
\\
 & TUBINO, Manoel José Gomes. Estudos brasileiros sobre o esporte: ênfase no esporte-educação.  Eduem, 2010.  163p.  ISBN: 9788576281771. }

\def\PUDautor{Adriano Barros}

\def\PUDdata{2017-07-20}

\def\PUDrevisao{1}

\def\PUDrevisor{Samuel Vieira Dias}

\def\PUDdatarevisao{2017-07-20}

\PUDcorpo
